\par\smallskip\noindent\textbf{完了通知だけを変えた追加診断．}
上の保存状態照合後，F/Bの２状態を新しい模擬環境で再現し，完了APIのなし／１回追加を固定４条件で比較した．送信済みの公開APIの返却と介入前13保存ファイルを全て照合し，既存Dと同じ空引数の通知だけを加えた．表4のとおりFは未達から成功へ変わり，その最終保存ファイルは元Dと全て一致した．Bは通知後も未達だった．非supervisorの保存ファイルは全条件で介入前後不変だった．
\begin{center}\small 表4\quad 別環境の再現における終了手続きの介入\par\smallskip
\begin{tabular}{lll}\hline 元状態 & 完了API追加 & 公式成否\\\hline
F & なし & 未達\\
F & １回追加 & 成功\\
B & なし & 未達\\
B & １回追加 & 未達\\
\hline\end{tabular}\end{center}
再現344と介入2の計346 API，独立評価４回，生成０で，全体23.927秒（評価8.541秒を含む）だった．これは既知１課題の事後診断であり，元のP55成否・未知費用は変更しない．自動で完了を宣言してよい条件や，新たな自律エージェントの成功率は検証していない．
